\documentclass[11pt]{article}

\usepackage{amsmath,amssymb}
\usepackage{booktabs}
\usepackage{hyperref}
\usepackage[margin=1in]{geometry}
\usepackage{longtable}
\usepackage{listings}
\usepackage{xcolor}

\lstset{
  language=Mathematica,
  basicstyle=\ttfamily\small,
  keywordstyle=\bfseries,
  commentstyle=\itshape\color{gray},
  breaklines=true,
  frame=single,
  columns=fullflexible,
  morekeywords={Module,Table,Do,If,Which,SetPrecision,FindRoot,NSolve,
    Association,OptionValue,Re,Im,Abs,Log,Exp,Sum,Product,Sort,
    Ordering,Length,Total,Get,Print,
    nLayerK2Num,nLayerDetNum,nLayerResidues,computePrefactor,
    computeFconv,computeHomogeneous}
}

\lstdefinestyle{shell}{
  language=bash,
  basicstyle=\ttfamily\small,
  frame=single,
  columns=fullflexible
}

\hypersetup{colorlinks=true,linkcolor=blue,citecolor=blue,urlcolor=blue}

\title{Constructing the Equivalent Homogeneous Rheology\\
  from a Multilayer Viscoelastic Model}
\author{}
\date{}

\begin{document}
\maketitle
\tableofcontents
\newpage

%======================================================================
\section{Overview}
\label{sec:overview}

This folder contains the modular Wolfram Language pipeline for
computing an equivalent homogeneous generalized Voigt rheology from an
arbitrary $N$-layer viscoelastic body.  The model is configured
entirely through \texttt{ModelInput.wl} and is not hardcoded to any
particular body.

The method follows Gevorgyan, Matsuyama \& Ragazzo (2023, MNRAS 523,
1822; hereafter GMR\,2023), which shows that the tidal Love number
$k_2(s)$ of any $N$-layer Maxwell body can be exactly reproduced by a
single homogeneous body with a generalized Voigt rheology.  The
layered-body formulation uses the propagator matrix notation of
Sabadini, Vermeersen \& Cambiotti (2016, hereafter SVC\,2016).

The pipeline proceeds in three stages:
\begin{enumerate}
  \item \textbf{Forward problem:} compute $k_2(s)$ as a rational
    function of the Laplace variable $s$ using the propagator
    $\mathbf{\Pi}_\ell$ (SVC\,2016, Eq.~2.9).
  \item \textbf{Modal decomposition:} find the secular poles
    $\{s_j\}$ and residues $\{r_j\}$ via grid search on the secular
    determinant $D(s)$ (SVC\,2016, Eq.~1.199), Brent refinement, and
    the limit formula.
  \item \textbf{Rheological inversion:} map the poles and residues to
    the generalized Voigt compliance $C(s) = P_2(s)/P_1(s)$
    (GMR\,2023, Eq.~6) and extract parameters
    $(\gamma, \alpha, \eta, \alpha_i, \eta_i)$,
    then identify dominant modes for a reduced model.
\end{enumerate}

%======================================================================
\section{File Inventory}
\label{sec:files}

\begin{longtable}{@{}lp{8.5cm}@{}}
\toprule
\textbf{File} & \textbf{Role} \\
\midrule
\endfirsthead
\toprule
\textbf{File} & \textbf{Role} \\
\midrule
\endhead
\texttt{ModelInput.wl} &
  User configuration: layer radii, densities, shear moduli,
  viscosities, tidal frequency, and numerical settings.  Edit this
  file to change the body model. \\[4pt]
\texttt{LayeredBodyConstants.wl} &
  Module~1: gravitational constant, gravity profile, Maxwell complex
  modulus, moment of inertia, mapping prefactor $\Theta$, conversion
  factor $F$, formatting. \\[4pt]
\texttt{YMatrix.wl} &
  Module~2: $6\times6$ fundamental matrix $\mathbf{Y}_\ell(r)$
  and propagator $\mathbf{\Pi}_\ell(r,r')$ for degree
  $\ell=2$ (SVC\,2016, Eq.~2.42).  Core--mantle boundary matrix
  $\mathbf{I}_C$ (Eq.~2.6). \\[4pt]
\texttt{NLayerJacobian.wl} &
  Module~5: $N$-layer propagator using the unified $6\times6$
  $\mathbf{\Pi}_\ell$ propagation through all layers (solid and
  liquid).  Liquid layers are handled via $\hat\mu \approx 10^{-20}$
  regularization.  The secular determinant $D(s)$
  (SVC\,2016, Eq.~1.199) is computed from the $3\times3$ projected
  matrix.  Provides \texttt{nLayerK2Num} (numerical $k_2(s)$),
  \texttt{nLayerDetNum} ($D(s)$), \texttt{nLayerResidues}
  (residues via limit formula), and perturbation utilities. \\[4pt]
\texttt{HomogeneousRheology.wl} &
  Module~6: \texttt{computeHomogeneous} --- the main function.  Calls
  the $N$-layer modules to find poles, compute residues, and extract
  all Voigt parameters.  Returns an \texttt{Association} with 16
  keys. \\[4pt]
\texttt{RunHomogeneous.wl} &
  Driver script.  Loads all modules and \texttt{ModelInput.wl},
  calls \texttt{computeHomogeneous}, identifies dominant modes, prints
  tables, and generates comparison plots. \\[4pt]
\texttt{VisualizeLayers.wl} &
  Standalone script that reads \texttt{ModelInput.wl} and draws a
  sector cross-section with numbered layers and a text listing of
  radii and thicknesses.  Uses layer names from
  \texttt{"layerNames"} if provided, otherwise labels generically.
  Exports \texttt{output/layer\_cross\_section.pdf}.  No pipeline
  modules required. \\
\bottomrule
\caption{Files in this folder.}
\label{tab:files}
\end{longtable}

\subsection{Module dependency graph}

\begin{verbatim}
RunHomogeneous.wl
  |-- ModelInput.wl
  |-- HomogeneousRheology`
  |     |-- LayeredBodyConstants`
  |     +-- NLayerJacobian`
  |           |-- LayeredBodyConstants`
  |           +-- YMatrix`

VisualizeLayers.wl
  +-- ModelInput.wl       (standalone, no modules)
\end{verbatim}

\subsection{Required load order}

\begin{lstlisting}
Get["LayeredBodyConstants.wl"];
Get["YMatrix.wl"];
Get["NLayerJacobian.wl"];
Get["HomogeneousRheology.wl"];
\end{lstlisting}

%======================================================================
\section{Quick Start}
\label{sec:quickstart}

\subsection{Prerequisites}

\begin{itemize}
  \item \textbf{Wolfram Engine} (free) or \textbf{Mathematica}
    (licensed), version 12.0 or later.
  \item \texttt{wolframscript} on your \texttt{PATH}.
\end{itemize}

\subsection{Running}

From the \texttt{homogeneous\_model/} directory:

\begin{lstlisting}[style=shell]
wolframscript -file RunHomogeneous.wl
\end{lstlisting}

This loads the model from \texttt{ModelInput.wl}, computes the
equivalent homogeneous Voigt body, identifies dominant relaxation
modes, prints all results, and generates four Love number comparison
plots.

\subsection{Changing the body model}

Edit \texttt{ModelInput.wl}.  All arrays must have the same length.
Set $\mu < 1$\,Pa for liquid layers.  Use exact rationals for maximum
precision.

%======================================================================
\section{Configuration: \texttt{ModelInput.wl}}
\label{sec:config}

\subsection{\texttt{\$ModelConfig} --- Physical model}

\begin{lstlisting}
$ModelConfig = <|
  "rk"   -> {r1, r2, ..., rN},   (* outer radii, m *)
  "rhok" -> {rho1, ..., rhoN},   (* densities, kg/m^3 *)
  "muk"  -> {mu1, ..., muN},     (* shear moduli, Pa *)
  "etak" -> {eta1, ..., etaN},   (* viscosities, Pa s *)
  "omega" -> 2 Pi / (T * 86400), (* tidal frequency, rad/s *)
  "layerNames" -> {"Layer 1", ..., "Layer N"}  (* optional *)
|>;
\end{lstlisting}

\begin{table}[h]
\centering
\begin{tabular}{@{}llp{7cm}@{}}
\toprule
\textbf{Key} & \textbf{Units} & \textbf{Description} \\
\midrule
\texttt{"rk"}    & m      & Outer radius of each layer, center to
  surface. \\
\texttt{"rhok"}  & kg/m$^3$ & Density of each layer. \\
\texttt{"muk"}   & Pa     & Shear modulus.  Set $\mu < 1$\,Pa for
  liquid layers. \\
\texttt{"etak"}  & Pa\,s  & Viscosity (any small value for liquid). \\
\texttt{"omega"} & rad/s  & Tidal forcing frequency. \\
\texttt{"layerNames"} & ---  & Optional list of $N$ strings for
  layer labels in \texttt{VisualizeLayers.wl}.  If omitted, layers are
  labelled generically as ``Layer~$i$ (solid/liquid)''. \\
\bottomrule
\end{tabular}
\caption{\texttt{\$ModelConfig} keys.}
\end{table}

\subsection{\texttt{\$AnalysisConfig} --- Numerical settings}

\begin{table}[h]
\centering
\begin{tabular}{@{}llp{6.5cm}@{}}
\toprule
\textbf{Key} & \textbf{Default} & \textbf{Description} \\
\midrule
\texttt{"workingPrecision"} & 50 & Digits of precision. \\
\texttt{"gridLogMin"} & $-20$ & $\log_{10}|s|$ lower bound for
  pole search. \\
\texttt{"gridLogMax"} & $-3$ & $\log_{10}|s|$ upper bound. \\
\texttt{"gridPoints"} & 2000 & Grid points for sign-change
  search. \\
\texttt{"omegaInterestMin"} & $10^{-10}$ & Lower bound of frequency
  interest range (rad/s). \\
\texttt{"omegaInterestMax"} & $10^{-5}$ & Upper bound (rad/s). \\
\texttt{"dominantThreshold"} & $1/50$ & Fractional threshold for
  dominant-mode selection (2\%). \\
\bottomrule
\end{tabular}
\caption{\texttt{\$AnalysisConfig} keys.}
\end{table}

%======================================================================
\section{Main Function: \texttt{computeHomogeneous}}
\label{sec:api}

\begin{lstlisting}
computeHomogeneous[rk, rhok, muk, etak, omega, opts]
\end{lstlisting}

\subsection{Algorithm}

\begin{enumerate}
  \item \textbf{Elastic Love number.} Evaluate $k_E =
    \operatorname{Re}[k_2(10^{10})]$.
  \item \textbf{Tidal $k_2$.} Evaluate $k_2(i\omega)$.
  \item \textbf{Pole-finding.} Grid search for zeros of $D(s)$
    (SVC\,2016, Eq.~1.199) on the negative real $s$-axis via
    \texttt{nLayerDetNum}, detect sign changes, refine with Brent.
  \item \textbf{Residues.} Limit formula
    $r_j = \delta\cdot k_2(s_j + \delta)$ with
    $\delta = 10^{-6}|s_j|$.
  \item \textbf{Fluid Love number.}
    $k_f = k_E + \sum(-r_j/s_j)$
    (GMR\,2023, Eq.~26).
  \item \textbf{Mapping constants.} $\Theta = 3\mathrm{I}_\circ G/R^5$
    (GMR\,2023, Eq.~2),
    $\gamma = \Theta/k_f$ (Eq.~4),
    $\alpha = \Theta/k_E - \gamma$,
    $F = 5\bar\rho R^2/38$.
  \item \textbf{Voigt inversion.} Build compliance
    $C(s) = P_2(s)/P_1(s)$ (Eq.~6), factor $s=0$ root from $P_1$
    for $\eta$, solve $P_2(s) = 0$ for Voigt poles $\{q_i\}$
    (Eq.~20), extract $(\alpha_i, \eta_i)$ from partial fractions
    of $C_v(s)$ (Eq.~23).
\end{enumerate}

\subsection{Return value}

Returns a Wolfram \texttt{Association} with 16 keys:

\begin{longtable}{@{}llp{6.5cm}@{}}
\toprule
\textbf{Key} & \textbf{Units} & \textbf{Description} \\
\midrule
\endfirsthead
\toprule
\textbf{Key} & \textbf{Units} & \textbf{Description} \\
\midrule
\endhead
\texttt{"kE"}        & dimless & Elastic Love number $k_E$. \\
\texttt{"kf"}        & dimless & Fluid Love number $k_f$. \\
\texttt{"k2Re"}      & dimless & $\operatorname{Re}[k_2(i\omega)]$. \\
\texttt{"k2Im"}      & dimless & $\operatorname{Im}[k_2(i\omega)]$. \\
\texttt{"gamma"}     & Pa & Gravitational rigidity $\gamma$
  (GMR\,2023, Eq.~4). \\
\texttt{"alpha"}     & Pa & Elastic rigidity $\alpha$. \\
\texttt{"eta"}       & Pa\,s & Isolated dashpot $\eta$. \\
\texttt{"prefactor"} & s$^{-2}$ & $\Theta = 3\mathrm{I}_\circ G/R^5$
  (GMR\,2023, Eq.~2). \\
\texttt{"Fconv"}     & Pa\,s$^2$ & Conversion factor $F$. \\
\texttt{"nPoles"}    & int & Number of secular poles $M$. \\
\texttt{"poles"}     & 1/s & Pole locations $\{s_j\}$. \\
\texttt{"residues"}  & dimless & Residues $\{r_j\}$. \\
\texttt{"nVoigt"}    & int & Number of Voigt elements $N_V$. \\
\texttt{"tauV"}      & s & Relaxation times $\tau_i = -1/p_i$. \\
\texttt{"alphaV"}    & Pa & Voigt spring constants $\alpha_i$. \\
\texttt{"etaV"}      & Pa\,s & Voigt dashpot viscosities $\eta_i$. \\
\bottomrule
\caption{Return value keys of \texttt{computeHomogeneous}.}
\end{longtable}

%======================================================================
\section{Driver Output: \texttt{RunHomogeneous.wl}}
\label{sec:output}

The driver prints seven output sections:

\begin{enumerate}
  \item \textbf{Love numbers:} $k_E$, $k_f$,
    $\operatorname{Re}[k_2]$, $\operatorname{Im}[k_2]$.
  \item \textbf{Rheological parameters:} $\gamma$, $\alpha$,
    $\eta$, pole/element counts.
  \item \textbf{Secular relaxation modes table:} all $M$ poles with
    $\tau_j$, residues, sorted by decreasing $|s_j|$.
  \item \textbf{Voigt elements table:} all $N_V$ elements with
    $\tau_i$, $\alpha_i$, $\eta_i$, sorted by decreasing $\tau$.
  \item \textbf{Dominant mode analysis:} fractional contribution of
    each element to $-\operatorname{Im}[J]$ over the interest range.
  \item \textbf{Reduced model summary:} $k_2$ from dominant-only
    model with relative errors.
  \item \textbf{Comparison plots:}
    $-\operatorname{Im}[k_2]$ and $\operatorname{Re}[k_2]$ vs
    $\omega$: $N$-layer vs full Voigt, $N$-layer vs reduced Voigt.
\end{enumerate}

%======================================================================
\section{Mathematical Pipeline}
\label{sec:math}

\subsection{From multilayer to $k_2(s)$}

Each solid layer has Maxwell rheology with complex shear modulus
(SVC\,2016, Correspondence Principle):
\begin{equation}
  \hat\mu_k(s) = \frac{\mu_k\,\nu_k\,s}{\mu_k + \nu_k\,s}\,,
\end{equation}
where $\mu_k$ is the shear modulus and $\nu_k$ the viscosity
(\texttt{"etak"} in the code).
Liquid layers use the regularization $\hat\mu \approx 10^{-20}$\,Pa,
which makes the $6\times6$ fundamental matrix $\mathbf{Y}_\ell(r)$
(SVC\,2016, Eq.~2.42) nearly rank-5 (the $S$ row vanishes).  All
layers are propagated through the full $6\times6$ propagator
$\mathbf{\Pi}_\ell(r, r') = \mathbf{Y}_\ell(r)\,
\mathbf{Y}_\ell^{-1}(r')$ (Eq.~2.9) chain from the core--mantle
boundary matrix $\mathbf{I}_C$ (Eq.~2.6); no special Clairaut
propagator is needed.  This yields $k_2(s)$ as a rational function
of $s$ for arbitrary layer configurations (any combination of solid
and liquid layers).

\subsection{Partial fraction decomposition}

The poles $s_j < 0$ of $k_2(s)$ are the secular relaxation modes, i.e.\
the roots of the secular determinant $D(s) = 0$ (SVC\,2016,
Eq.~1.200).  The partial-fraction form is (GMR\,2023, Eq.~26):
\begin{equation}
  k_2(s) = k_E + \sum_{j=1}^{M} \frac{r_j}{s - s_j}\,,
\end{equation}
where $r_j$ are the residues and $\tau_j = -1/s_j$ the relaxation
times of the normal modes.

\subsection{Compliance inversion}

The compliance is written as a ratio of two polynomials
(GMR\,2023, Eq.~6):
\begin{equation}
  C(s) = \frac{P_2(s)}{P_1(s)}
    = \frac{k_2(s)}{\Theta - \gamma\,k_2(s)}\,,
\end{equation}
where $\Theta = 3\mathrm{I}_\circ G/R^5$ (Eq.~2) and
$\gamma = \Theta/k_f$ is the gravitational rigidity (Eq.~4).
The viscous part of the compliance decomposes into partial fractions
as the generalized Voigt form (Eq.~23):
\begin{equation}
  C_v(s) = \frac{\alpha^2}{(\gamma+\alpha)^2}
    \biggl(\frac{1}{\eta} + \sum_{i=1}^{n}\frac{1}{\eta_i}\biggr)
    \biggl(\frac{A_1}{s - s_1} + \cdots
    + \frac{A_{n+1}}{s - s_{n+1}}\biggr),
\end{equation}
from which the full compliance is (Eq.~14):
\begin{equation}
  C(s) = \frac{1}{\alpha + \gamma} + C_v(s)\,.
\end{equation}
The isolated dashpot $\eta$ and the Kelvin--Voigt element parameters
$\alpha_i$, $\eta_i$ (with characteristic times
$\omega_i^{-1} = \eta_i/\alpha_i$, Eq.~18)
are extracted from the partial-fraction coefficients.
Physical (Pa) values: multiply by $F = 5\bar\rho R^2/38$.

\subsection{Dominant mode selection}

Each Voigt element's maximum fractional contribution to
$-\operatorname{Im}[C(i\omega)]$ is computed over
$[\omega_{\min}, \omega_{\max}]$.  Elements exceeding the threshold
are retained in the reduced model.

%======================================================================
\section{Customising for a Different Body}
\label{sec:custom}

\begin{enumerate}
  \item Edit \texttt{ModelInput.wl}: set \texttt{"rk"},
    \texttt{"rhok"}, \texttt{"muk"}, \texttt{"etak"} for your body.
    All arrays must have the same length $N$.
  \item Set \texttt{"omega"} to the body's tidal forcing frequency.
  \item Adjust \texttt{"gridLogMin"}/\texttt{"gridLogMax"} if
    relaxation times span a different range than the default.
  \item Adjust \texttt{"omegaInterestMin"}/\texttt{"omegaInterestMax"}
    to the frequency range relevant for your problem.
  \item Run: \texttt{wolframscript -file RunHomogeneous.wl}.
\end{enumerate}

The pipeline auto-detects the number of layers, identifies liquid
layers ($\mu < 1$\,Pa), and adjusts the propagator chain accordingly.
No code changes are needed.

%======================================================================
\section{Test Cases}
\label{sec:testcases}

The \texttt{test\_cases/} folder contains input files for three
tidally-heated bodies beyond the default 5-layer Moon model.  These
validate the pipeline for different layer counts and liquid-layer
patterns.

\begin{description}
  \item[\texttt{ModelInputEuropa.wl}] --- 4~layers (liquid iron core,
    silicate mantle, subsurface ocean, ice shell).  Two liquid layers.
    Tidal period 3.5512\,days.  Based on Wahr et al.\ (2006),
    Hussmann \& Spohn (2004).
  \item[\texttt{ModelInputEnceladus.wl}] --- 3~layers (silicate core,
    subsurface ocean, ice shell).  One liquid layer.  Tidal period
    1.3702\,days.  Based on Cadek et al.\ (2016), Thomas et al.\
    (2016).
  \item[\texttt{ModelInputIo.wl}] --- 4~layers (liquid Fe-FeS core,
    mantle, asthenosphere, lithosphere).  One liquid layer, contrasting
    solid viscosities.  Tidal period 1.7691\,days.  Based on Segatz
    et al.\ (1988), Bierson \& Nimmo (2016).
\end{description}

\subsection{Running a test case}

Copy the desired input file over \texttt{ModelInput.wl} and run the
pipeline:

\begin{lstlisting}[style=shell]
cd homogeneous_model/
cp test_cases/ModelInputEuropa.wl ModelInput.wl
wolframscript -file RunHomogeneous.wl
wolframscript -file VisualizeLayers.wl
\end{lstlisting}

Repeat for Enceladus or Io by substituting the file name.  Remember to
restore the original \texttt{ModelInput.wl} afterwards if needed.

%======================================================================
\section{Layer Visualization: \texttt{VisualizeLayers.wl}}
\label{sec:vis}

Standalone script that generates a sector cross-section of the layered
body.  Works for any number of layers.

\begin{lstlisting}[style=shell]
wolframscript -file VisualizeLayers.wl
\end{lstlisting}

\paragraph{Input.}
Reads \texttt{ModelInput.wl} (\texttt{"rk"}, \texttt{"muk"}, and
optionally \texttt{"layerNames"}).

\paragraph{Liquid detection.}
$\mu < 1$\,Pa $\Rightarrow$ liquid.

\paragraph{Colours.}
Layers are coloured with the \texttt{LightTemperatureMap} palette: warm
(orange) at the center, cool (blue) at the surface.  The range is
compressed to $[0.2, 0.8]$ to avoid extreme saturation.

\paragraph{Layer names.}
If \texttt{"layerNames"} is provided in \texttt{\$ModelConfig}, those
names are used directly.  Otherwise, layers are labelled generically as
``Layer~$i$ (solid)'' or ``Layer~$i$ (liquid)'' based on the liquid
detection.  The pre-configured test cases include descriptive names
(e.g.\ ``Inner core'', ``Silicate mantle'', ``Ice shell'').

\paragraph{Thin-layer exaggeration.}
If any layer occupies less than 3\% of the surface radius, its visual
thickness is boosted to the minimum and the deficit is redistributed
proportionally from thicker layers.  The title is updated to note
``(visual thickness exaggerated)'' when this occurs; otherwise it reads
``(to scale)''.

\paragraph{Layout.}
The cross-section sector is on the left with layer numbers placed along
the left radial edge at each layer's midpoint, and ``0'' at the body
centre.  On the right, a text listing shows each layer as:
\texttt{i --- Name~~thickness~km (radius~km)}.

\paragraph{Output.}
\texttt{output/layer\_cross\_section.pdf} --- sector cross-section with
numbered layers on the left, text listing of names, thicknesses, and
outer radii on the right.

%======================================================================
\section{Troubleshooting}
\label{sec:trouble}

\begin{description}
  \item[Missing poles:] Widen \texttt{"gridLogMin"}/\texttt{"gridLogMax"}
    or increase \texttt{"gridPoints"}.
  \item[Curves don't overlap in full Voigt plot:] Increase
    \texttt{"workingPrecision"} (try 80 or 100).
  \item[\texttt{FindRoot} warnings:] Usually harmless.  If poles look
    wrong, increase grid density.
  \item[\texttt{NSolve::precw}:] The Voigt pole solver needed extra
    precision; handled automatically at doubled precision.
\end{description}

\bigskip
\noindent\textbf{References.}

\medskip
\noindent
\textit{Gevorgyan, Matsuyama \& Ragazzo (2023), MNRAS 523, 1822.}

\medskip
\noindent
\textit{Sabadini, Vermeersen \& Cambiotti (2016), Global Dynamics of
the Earth: Applications of Viscoelastic Relaxation Theory to
Solid-Earth and Planetary Geophysics, 2nd ed., Springer.}

\end{document}
